% SPDX-License-Identifier: CC-BY-SA-4.0
% Copyright (c) 2025 Luis Alonso
%
% This file is part of luis_alonso/Notas-Japones.
%
% Licensed under Creative Commons Attribution-ShareAlike 4.0 International (CC BY-SA 4.0).
% You may share and adapt this material for any purpose,
% as long as you provide attribution and distribute adaptations under the same license.
%
% License: https://creativecommons.org/licenses/by-sa/4.0/
% Source:  https://git.lalonso.com/luis_alonso/Notas-Japones
%
% Note: Third-party content (if any) is excluded from this license and is marked where used.

\documentclass[a4paper,lualatex,ja=standard,base=10pt]{bxjsarticle}

\title{\ruby{漢|字}{かん|じ} LV5}
\author{アロンゾ　ルイス}

\setmainjfont{Noto Serif CJK JP}
\usepackage{luatexja-fontspec}
\usepackage{kanjicard} % Custom Kanji card
\usepackage{fancyhdr}

% --- Page style ---
\pagestyle{fancy}
\fancyhf{} % clear header/footer

% Footer: name + copyright on left, page number on right
\fancyfoot[L]{Luis Alonso\ \copyright\ \the\year}
\fancyfoot[R]{\thepage}

% Optional: remove header rule and keep a footer rule
\renewcommand{\headrulewidth}{0pt}
\renewcommand{\footrulewidth}{0.4pt}

% Increase space reserved for header area
\setlength{\headheight}{0pt}
\setlength{\headsep}{6pt} % distance from header to text

\fancypagestyle{plain}{%
  \fancyhf{}
  \fancyfoot[L]{Luis Alonso\ \copyright\ \the\year}
  \fancyfoot[R]{\thepage}
  \renewcommand{\headrulewidth}{0pt}
  \renewcommand{\footrulewidth}{0.4pt}
}

\begin{document}
\setcounter{page}{1} % start page number here

\raggedbottom
\maketitle{}
\thispagestyle{fancy}

\begin{kanjicard}{時}{とき}{ジ}
  \KExample{時}{とき}{Tiempo}
  \KExample{四時}{よ・じ}{4:00}
  \KExample{一時間}{いち・じ・かん}{Una hora}
  \KExample{時計}{と・けい}{Reloj}
\end{kanjicard}

\begin{kanjicard}{間}{あいだ/ま}{カン}
  \KExample{間}{あいだ}{Entre}
  \KExample{一年間}{いち・ねん・かん}{Un año}
  \KExample{間に合う}{ま・に・あ・う}{Llegar a tiempo}
  \KExample{時間}{じ・かん}{Tiempo}
\end{kanjicard}

\begin{kanjicard}{場}{ば}{ジョウ}
  \KExample{場所}{ばしょ}{Lugar}
  \KExample{場合}{ば・あい}{Caso}
  \KExample{広場}{ひろ・ば}{Plaza}
  \KExample{サッカー場}{サッカー・じょう}{Cancha de futbol}
\end{kanjicard}

\begin{kanjicard}{所}{ところ}{ショ/ジョ}
  \KExample{所}{ところ}{Lugar}
  \KExample{研究所}{けん・きゅう・じょ}{Instituto de investigación}
  \KExample{台所}{だい・どころ}{Cocina}
  \KExample{住所}{じゅう・しょ}{Dirección}
\end{kanjicard}

\begin{kanjicard}{駅}{}{エキ}
  \KExample{駅}{えき}{Estación}
  \KExample{駅員}{えき・いん}{Empleado de estación}
  \KExample{駅前}{えき・まえ}{Frente a la estación}
\end{kanjicard}

\begin{kanjicard}{出}{で・る/で・かける/だ・す}{ジュッ/ジュツ}
  \KExample{出かける}{で・かける}{Salir}
  \KExample{出ます}{で・ます}{Salir}
  \KExample{出します}{だ・します}{Entregar/Sacar}
  \KExample{出発}{しゅっ・ぱつ}{Partir}
\end{kanjicard}

\begin{kanjicard}{待}{ま・つ/まち}{タイ}
  \KExample{待つ}{ま・つ}{Esperar}
  \KExample{招待する}{じょう・たい・する}{Invitar}
  \KExample{待合室}{まち・あい・しつ}{Sala de espera}
\end{kanjicard}

\begin{kanjicard}{止}{と・まる/と・める}{シ}
  \KExample{止まる}{と・まる}{Detenerse} %が
  \KExample{止める}{と・める}{Detener} %を
  \KExample{中止}{ちゅう・し}{Cancelar/Suspender}
  \KExample{通行止め}{つう・こう・ど・め}{Calle cerrada}
\end{kanjicard}

\begin{kanjicard}{事}{こと}{ジ}
  \KExample{事}{こと}{Hecho, asunto, cosa}
  \KExample{事故}{じ・こ}{Accidente}
  \KExample{食事}{しょく・じ}{Comida}
  \KExample{用事}{よう・じ}{Compromiso, pendiente}
\end{kanjicard}

\begin{kanjicard}{仕}{つか・える}{シ}
  \KExample{仕事}{し・ごと}{Trabajo}
  \KExample{仕える}{つか・える}{Servir}
  \KExample{仕方}{し・かた}{Forma de hacer algo}
\end{kanjicard}

\begin{kanjicard}{前}{まえ}{ゼン}
  \KExample{前}{まえ}{Antes, adelante}
  \KExample{-年前}{ねん・まえ}{Años antes}
  \KExample{前半}{ぜん・はん}{Primera mitad}
\end{kanjicard}

\begin{kanjicard}{後}{あとうし・ろ}{ゴ/コウ}
  \KExample{後}{あと}{Después}
  \KExample{-年後}{ねん・ご}{Años después}
  \KExample{後ろ}{うしろ}{Atrás}
  \KExample{後半}{こう・はん}{Segunda mitad}
\end{kanjicard}

\begin{kanjicard}{朝}{あさ}{チョウ}
  \KExample{朝}{あさ}{Mañana}
  \KExample{朝食}{ちょうしょく}{Desayuno}
  \KExample{朝日}{あさ・ひ}{Sol de la mañana}
\end{kanjicard}

\begin{kanjicard}{昼}{ひる}{チュウ}
  \KExample{昼}{ひる}{Tarde}
  \KExample{昼食}{ちゅう・しょく}{Comida/Almuerzo}
  \KExample{昼寝}{ひる・ね}{Siesta}
  \KExample{昼休み}{ひる・やすみ}{Receso del almuerzo}
\end{kanjicard}

\begin{kanjicard}{夕}{ゆう}{}
  \KExample{夕方}{ゆう・がた}{Tarde-noche}
  \KExample{夕食}{ゆう・しょく}{Merienda/Cena}
  \KExample{七夕}{たな・ばた}{Festival de las estrellas}
\end{kanjicard}

\begin{kanjicard}{夜}{よる/よ}{や}
  \KExample{夜}{よる}{Noche}
  \KExample{今夜}{こん・や}{Esta noche}
  \KExample{夜食}{や・しょく}{Bocadillo de media noche}
  \KExample{夜空}{よ・ぞら}{Cielo nocturno}
\end{kanjicard}

\begin{kanjicard}{乗}{の・る}{ジョウ}
  \KExample{乗る}{のる}{Subir}
  \KExample{乗車する}{じょう・しゃ・する}{Subir a un medio de transp.}
  \KExample{乗客}{じょう・きゃく}{Pasajero}
\end{kanjicard}
% じょうしゃけん 乗車券 boleto para transporte

\begin{kanjicard}{学}{まな・ぶ}{ガク}
  \KExample{学ぶ}{まな・ぶ}{Aprender}
  \KExample{大学}{だい・がく}{Universidad}
  \KExample{入学}{にゅう・がく}{Entrar a la escuela}
\end{kanjicard}

\begin{kanjicard}{校}{}{コウ}
  \KExample{学校}{がっ・こう}{Escuela}
  \KExample{小学校}{しょう・がっ・こう}{Primaria}
  \KExample{中学校}{ちゅう・がっ・こう}{Secundaria}
  \KExample{高校}{こう・こう}{Preparatoria}
\end{kanjicard}

\begin{kanjicard}{先}{さき/ま・ず}{セン}
  \KExample{先}{さき}{Adelante, antes de, primero}
  \KExample{先月}{せん・げつ}{Mes pasado}
  \KExample{先日}{せん・じつ}{El otro día, hace unos días}
\end{kanjicard}

\begin{kanjicard}{生}{い・きる/う・まれる}{セイ/ショウ}
  \KExample{生きる}{い・きる}{Vivir}
  \KExample{生まれる}{う・まれる}{Nacer}
  \KExample{先生}{せん・せい}{Maestro}
  \KExample{学生}{がく・せい}{Estudiante}
  \KExample{-年生}{ねん・せい}{Estudiante de tal año}
\end{kanjicard}

\begin{kanjicard}{勉}{}{ベン}
  \KExample{勉強}{べん・きょう}{Estudio}
\end{kanjicard}

\begin{kanjicard}{強}{つよ・い}{キョウ/ゴウ}
  \KExample{強い}{つよ・い}{Fuerte}
  \KExample{勉強}{べん・きょう}{Estudio}
  \KExample{強盗}{ごう・とう}{Robo, atraco}
\end{kanjicard}

% Segunda parte

\begin{kanjicard}{文}{}{ブン/モン}
  \KExample{文}{ぶん}{Oración}
  \KExample{注文}{ちゅう・もん}{Orden}
  \KExample{文学}{ぶん・がく}{Literatura}
\end{kanjicard}

\begin{kanjicard}{化}{}{カ/ケ}
  \KExample{文化}{ぶん・か}{Cultura, civilización}
  \KExample{化学}{か・がく}{Química}
  \KExample{変化}{へん・か}{Cambio}
\end{kanjicard}

\begin{kanjicard}{音}{おと}{オン/イン}
  \KExample{音}{おと}{Sonido, ruido}
  \KExample{発音}{はつ・おん}{Pronunciación}
  \KExample{録音}{ろく・おん}{Grabación}
\end{kanjicard}

\begin{kanjicard}{楽}{たの・しい}{ガク/ラク/ゴウ}
  \KExample{楽しい}{たの・しい}{Divertido}
  \KExample{楽}{らく}{Cómodo, sencillo}
  \KExample{音楽}{おん・がく}{Música}
\end{kanjicard}

\begin{kanjicard}{旅}{たび}{リョ}
  \KExample{旅}{たび}{Viaje (expresa aventura)}
  \KExample{旅行}{りょ・こう}{Viaje, viajar}
  \KExample{旅館}{りょ・かん}{Hotel (estilo japonés)}
\end{kanjicard}

\begin{kanjicard}{留}{とど・まる}{リュウ/ル}
  \KExample{留まる}{とど・まる}{Parar, quedarse}
  \KExample{留学}{りゅう・がく}{Estudios en el extranjero}
  \KExample{留守}{る・す}{No estar en casa}
\end{kanjicard}

\begin{kanjicard}{友}{とも}{ユウ}
  \KExample{友達}{とも・だち}{Amigo}
  \KExample{親友}{しん・ゆう}{Amigo íntimo}
\end{kanjicard}

\begin{kanjicard}{週}{}{シュウ}
  \KExample{週}{しゅう}{Semana}
  \KExample{週末}{しゅう・まつ}{Fin de semana}
  \KExample{先週}{せん・しゅう}{Semana pasada}
\end{kanjicard}

\begin{kanjicard}{回}{まわ・る/まわ・す}{カイ/エ}
  \KExample{回る}{まわ・る}{Girar}
  \KExample{2回}{に・かい}{Dos veces}
  \KExample{何回}{なん・かい}{¿Cuántas veces?}
\end{kanjicard}

% Leccion 11
\begin{kanjicard}{物}{もの}{ブツ}
  \KExample{物}{もの}{Cosa, objeto}
  \KExample{食べ物}{た・べ・もの}{Comida}
  \KExample{動物}{どう・ぶつ}{Animales}
\end{kanjicard}

\begin{kanjicard}{茶}{}{チャ/サ}
  \KExample{お茶}{お・ちゃ}{Té verde}
  \KExample{茶道}{さ・どう}{Ceremonia del té}
  \KExample{茶色}{ちゃ・いろ}{Color café}
\end{kanjicard}

\begin{kanjicard}{酒}{さけ}{シュ}
  \KExample{酒}{さけ}{Sake}
  \KExample{日本酒}{に・ほん・しゅ}{Vino de arroz japonés}
  \KExample{飲酒}{いん・しゅ}{Beber alcohol}
\end{kanjicard}

\begin{kanjicard}{作}{つく・る}{サク/サ}
  \KExample{作る}{つく・る}{Hacer, crear, fabricar}
  \KExample{作文}{さく・ぶん}{Escrito, composición}
  \KExample{作家}{さっ・か}{Autor}
\end{kanjicard}

\begin{kanjicard}{持}{も・つ}{}
  \KExample{持つ}{もつ}{Poseer, llevar, cargar}
  \KExample{持っていく}{も・っていく}{Llevar}
  \KExample{気持ち}{き・も・ち}{Sentimiento, sensación}
\end{kanjicard}

\begin{kanjicard}{願}{ねが・い}{ガン}
  \KExample{願い}{ねが・い}{Deseo}
  \KExample{お願いします}{お・ねが・いします}{Por favor}
  \KExample{願書}{がん・しょ}{Solicitud (por escrito)}
\end{kanjicard}

% Lección 12

\begin{kanjicard}{料}{}{リョウ}
  \KExample{無料}{む・りょう}{Gratis}
  \KExample{材料}{ざい・りょう}{Ingredientes, material}
  \KExample{給料}{きゅう・りょう}{Sueldo, salario}
\end{kanjicard}

\begin{kanjicard}{理}{}{リ}
  \KExample{料理}{りょう・り}{Cocina, platillo}
  \KExample{理解}{り・かい}{Comprensión}
  \KExample{理由}{り・ゆう}{Razón, motivo}
\end{kanjicard}

\begin{kanjicard}{味}{あじ}{ミ}
  \KExample{味}{あじ}{Sabor, gusto}
  \KExample{意味}{い・み}{Significado}
  \KExample{趣味}{しゅ・み}{Pasatiempo}
\end{kanjicard}

\begin{kanjicard}{色}{いろ}{ショク/シキ}
  \KExample{色}{いろ}{Color}
  \KExample{色々}{いろ・いろ}{Diversos, varios}
  \KExample{景色}{け・しき}{Paisaje}
  \KExample{三色}{さん・しょく}{Tres colores}
\end{kanjicard}

\begin{kanjicard}{野}{の}{ヤ/ショ}
  \KExample{野原}{の・はら}{Campo}
  \KExample{野球}{や・きゅう}{Beisbol}
\end{kanjicard}

\begin{kanjicard}{菜}{}{サイ}
  \KExample{野菜}{や・さい}{Vegetales}
  \KExample{菜園}{さい・えん}{Huerto}
\end{kanjicard}

\begin{kanjicard}{少}{すく・ない}{ショウ}
  \KExample{少ない}{すく・ない}{Poca cantidad}
  \KExample{少し}{すこ・し}{Un poco}
  \KExample{少々}{しょう・しょう}{Un poco}
\end{kanjicard}

\begin{kanjicard}{入}{はい・る/い・れる}{ニュウ/ジュ}
  \KExample{入る}{はい・る}{Entrar}
  \KExample{入れる}{い・れる}{Meter}
  \KExample{入学}{にゅう・がく}{Ingresar a la escuela}
\end{kanjicard}

\end{document}

